% Fuentes y límites de redacción: MAPA_FUENTES.md, capítulo 1.
\chapter{Introducción}
\label{ch:introduccion}

\section{Contexto y motivación}
La incorporación de instalaciones renovables al sistema eléctrico requiere,
además de su diseño y construcción, una tramitación administrativa. Conocer
esa tramitación ayuda a distinguir una propuesta en estudio de un proyecto
que ha recibido una autorización. Este seguimiento resulta de interés en el
contexto de la transición hacia fuentes de energía renovables.
\citapendiente{Incorporar una fuente institucional sobre transición energética
y planificación de renovables en España; no se incluyen cifras de despliegue
sin verificar.}

Una misma planta puede aparecer en distintas publicaciones oficiales: un
anuncio de información pública, una decisión ambiental o una autorización,
entre otras. Cada publicación explica una parte del procedimiento, pero no
constituye necesariamente una historia completa del proyecto. Además, una
resolución puede recordar actuaciones anteriores para justificar la decisión
que adopta.

Este trabajo utiliza el Boletín Oficial del Estado (BOE) como fuente de
información. El BOE permite consultar las publicaciones originales, aunque
seguir una planta exige relacionar documentos, reconocer sus denominaciones
y distinguir qué se publica ahora de lo que se cita como antecedente. Una
herramienta que reúna esas publicaciones puede facilitar la consulta, siempre
que conserve la relación entre cada dato y el texto que lo respalda.

\section{Problema abordado}
La pregunta que guía el trabajo es: \emph{¿cómo reconstruir la evolución
administrativa observada de un proyecto de generación eléctrica a partir de
la evidencia publicada en el BOE?}

El problema combina dos tareas. Primero, es necesario convertir documentos
escritos en lenguaje natural en información organizada: qué planta se
menciona, qué actuación administrativa se publica y a qué territorio se
refiere. Después, hay que determinar qué menciones de distintos documentos
corresponden al mismo proyecto. Compartir promotor, provincia o
infraestructura no basta para afirmar que dos plantas son una sola.

En esta memoria se denomina \emph{actuación administrativa} al trámite o
decisión que una publicación vincula a una planta o a un componente asociado.
La \emph{evidencia} es el pasaje del documento que permite comprobar esa
información. La reconstrucción resultante describe lo observado en las
publicaciones incluidas; no permite asegurar que se haya recuperado todo el
expediente administrativo.

\section{Objetivo general}
Desarrollar y evaluar un sistema reproducible que extraiga información de
publicaciones del BOE, relacione menciones de plantas de generación eléctrica
y permita consultar su evolución administrativa con trazabilidad hasta la
fuente documental. Un sistema es reproducible, en el sentido utilizado aquí,
cuando conserva los datos, las reglas y las versiones necesarias para
verificar sus resultados y repetir las transformaciones deterministas, es
decir, las que aplican reglas fijas a las mismas entradas.

\section{Objetivos específicos}
\begin{enumerate}
\item Obtener publicaciones candidatas del BOE mediante criterios explícitos
y conservar su identidad, fecha y procedencia.
\item Extraer plantas, actuaciones y territorios con una estructura definida
y evidencia documental.
\item Validar la información, conservar los casos inciertos para revisión
y aplicar decisiones humanas de forma trazable.
\item Relacionar menciones de plantas mediante reglas conservadoras y
vincular las referencias territoriales con unidades administrativas oficiales.
\item Preparar datos consolidados para una aplicación de consulta de solo
lectura, con acceso a proyectos, publicaciones y cronologías.
\item Evaluar la extracción y su procesamiento documental con un conjunto
independiente anotado por una persona, mediante reglas fijadas antes de
observar las predicciones del sistema.
\end{enumerate}

Estos objetivos separan la construcción del sistema de su evaluación
experimental. La evaluación independiente se centra en la salida documental
de extracción, canonicalización y validación. La canonicalización consiste
en aplicar reglas comunes para obtener una representación consistente de
esa salida. La agrupación posterior de proyectos y la aplicación cuentan con
validaciones técnicas, pero no reciben una medida de calidad independiente
a partir de este experimento.

\section{Alcance}
El objeto de seguimiento son plantas de generación eléctrica con nombre. Las
instalaciones de almacenamiento o evacuación pueden asociarse a una planta
cuando la publicación establece esa relación, pero no crean por sí solas un
proyecto de generación. El nombre del proyecto actúa como referencia para
reunir sus menciones; las reglas concretas se explican en el
capítulo~\ref{ch:metodologia}.

El corpus final, es decir, el conjunto de documentos utilizado para construir
el producto, combina publicaciones de una ventana reciente con una búsqueda
retrospectiva conservadora desde 2022. No es un censo exhaustivo de proyectos
renovables en España. Tampoco incluye sistemáticamente boletines autonómicos
o provinciales, expedientes completos ni fuentes de operación de las plantas.
El capítulo~\ref{ch:datos} precisa sus periodos y unidades.

La herramienta presenta información publicada; no certifica el estado
jurídico completo de un proyecto ni predice su siguiente autorización.
Los territorios reflejan asociaciones documentales y no coordenadas exactas
de las instalaciones. No se estima producción energética y el producto final
no incorpora indicadores consolidados de promotor o potencia. La operación
diaria automática, la administración de usuarios y una interfaz de aprobación
de correcciones quedan fuera del trabajo entregable.

La extracción utiliza un modelo de inteligencia artificial capaz de procesar
texto y generar respuestas, denominado modelo de lenguaje de gran tamaño o
\emph{LLM}. Su salida puede contener errores y requiere validación. Repetir
una consulta al proveedor no garantiza una respuesta idéntica; por ello se
conservan las predicciones originales utilizadas en la evaluación.

\section{Estructura de la memoria}
El capítulo~\ref{ch:contexto} introduce los principales trámites y el BOE.
El capítulo~\ref{ch:datos} explica cómo se obtuvieron y seleccionaron los
documentos. El capítulo~\ref{ch:metodologia} desarrolla la metodología de
extracción, validación y consolidación, y el capítulo~\ref{ch:aplicacion}
describe la aplicación de consulta.

El capítulo~\ref{ch:evaluacion} presenta el diseño experimental y define las
medidas necesarias para interpretar el capítulo~\ref{ch:resultados}. Después,
el capítulo~\ref{ch:discusion} analiza el significado y los límites de los
resultados. El capítulo~\ref{ch:conclusiones} responde a los objetivos y
organiza el trabajo futuro. Los anexos reservan el detalle de reproducibilidad
y del contrato documental que no resulta necesario para seguir la narración.
